\section{Signatures and terms}
\label{sec:fol-syntax-semantics-01}

Fix the countably infinite set
\(\Var\defeq\set{x_n:n\in\N}\) of object-language variables.  The letters
\(x,y,z\) denote members of \(\Var\), while \(r,t,u\) are metavariables
ranging over terms.  The letter \(s\) is reserved for variable assignments.

\begin{definition}[First-order signature]
A \emph{first-order signature} is a tuple
\[
 \Lang=\tuple{\mathcal F,\mathcal R,\arity_{\mathcal F},
                    \arity_{\mathcal R}},
\]
where \(\mathcal F\) and \(\mathcal R\) are sets of function and relation
symbols and
\[
 \arity_{\mathcal F}:\mathcal F\to\N,
 \qquad
 \arity_{\mathcal R}:\mathcal R\to\N
\]
assign each symbol its arity.  A function symbol of arity zero is a
\emph{constant symbol}.  The equality sign \(=\) is logical and is not a
member of \(\mathcal R\).  The sets \(\mathcal F\), \(\mathcal R\),
\(\Var\), and the set of reserved logical and punctuation tokens
\[
 \{=,\bottom,\neg,\land,\lor,\to,\forall,\exists,
   \texttt{(},\texttt{)},\texttt{,}\}
\]
are pairwise disjoint.  Consequently no signature symbol can also be a
variable, logical symbol, parenthesis, or comma.
\end{definition}

The cardinality of the signature is
\(\abs{\Lang}\defeq\abs{\mathcal F\cup\mathcal R}\).  Arity data add no
new symbols.

The signature is a syntactic object.  It assigns arities but assigns no
functions, relations, or elements.  A structure supplies those semantic
objects later.

\begin{definition}[Term]
For a signature \(\Lang\), the set \(\Term(\Lang)\) of \emph{terms} is the
least set satisfying:
\begin{enumerate}
  \item every \(x\in\Var\) belongs to \(\Term(\Lang)\); and
  \item if \(f\in\mathcal F\) has arity \(n\) and
        \(t_1,\ldots,t_n\in\Term(\Lang)\), then
        \(f(t_1,\ldots,t_n)\in\Term(\Lang)\).
\end{enumerate}
For \(n=0\), the expression is written \(f\), not \(f()\).  No other string
is a term.
\end{definition}

\begin{example}
Let \(\Lang_{\mathrm{ar}}\) have function symbols
\(\set{0,S,+,\cdot}\) of arities \(0,1,2,2\), respectively, and relation
symbol \(<\) of arity \(2\).  The strings
\[
 x_0,\qquad S(0),\qquad +(S(0),x_2),\qquad
 \cdot(+(x_0,x_1),S(x_2))
\]
are terms.  The string \(<(x_0,x_1)\) is not a term because \(<\) is a
relation symbol.  Conventional infix displays such as \(x_0+S(0)\) are
typographical abbreviations for the prefix grammar.
\end{example}

